%------------------------------------------------------------------------------%
\documentclass[fleqn,utf8,aspectratio=169,12pt]{beamer}

\usepackage{calculo_varias_variaveis-1}

\title{Derivação Implícita}

%------------------------------------------------------------------------------%
\begin{document}

\addtitle

%------------------------------------------------------------------------------%
\section{Derivação Implícita}

%------------------------------------------------------------------------------%
\begin{frame}{Derivação Implícita}
  Suponha que $F(x, y)$ seja diferenciável e que a equação
  \[
    F(x, y) = 0
  \]
  defina $y$ como função de $x$, encontre \(\frac{dy}{dx}\)
\end{frame}

%------------------------------------------------------------------------------%
\begin{proofframe}{Derivação Implícita -- Justificativa}
\begin{columns}
\column{0.4\linewidth}
\begin{align*}
  \onslide<+->{w(x)          = F\big(x, y(x)\big)                              & = 0 \\[4mm]}
  \onslide<+->{\frac{dw}{dx} = \D{F}{x} \frac{dx}{dx} + \D{F}{y} \frac{dy}{dx} & = 0}
\end{align*}
\column{0.6\linewidth}
\begin{align*}
  \onslide<+->{F_x \frac{dx}{dx} + F_y \frac{dy}{dx} & = 0      \\[2mm]}
  \onslide<+->{F_x + F_y \frac{dy}{dx}               & = 0      \\[2mm]}
  \onslide<+->{F_y \frac{dy}{dx}                     & = - F_x  \\[2mm]}
  \onslide<+->{\frac{dy}{dx}                         & = - \frac{F_x}{F_y}
               \qquad \text{desde que } F_y \neq 0}
\end{align*}
\end{columns}
\end{proofframe}

%------------------------------------------------------------------------------%
\begin{frame}{Derivação Implícita -- Fórmula}
  Suponha que $F(x, y)$ seja diferenciável e que a equação
  \[
    F(x, y)=0
  \]
  defina $y$ como função de $x$, \(y=y(x)\), então
  \[
    \frac{dy}{dx} = -\frac{F_x}{F_y}
  \]
  desde que \(F_y \neq 0\)
\end{frame}

%------------------------------------------------------------------------------%
\include{G-02-Derivacao_implicita-ex1}

%------------------------------------------------------------------------------%
\begin{frame}{Derivação Implícita -- Três Variáveis}
  Suponha que $F(x, y, z)$ seja diferenciável e que a equação
  \[
    F(x, y, z) = 0
  \]
  defina $z$ como função de $(x, y)$, isto é, $z=f(x, y)$,
  determine
  \(\D{z}{x}\) e
  \(\D{z}{y}\)
\end{frame}

%------------------------------------------------------------------------------%
\begin{proofframe}{Derivação Implícita -- Justificativa}
\begin{columns}
  \column{0.5\linewidth}
  \begin{align*}
    \onslide<+->{F\big(x, y, z(x, y)\big)                               & = 0 \\[2mm]}
    \onslide<+->{\D{} {x\vphantom{y}}\left[F\big(x, y, z(x, y)\big)\right]          & = 0 \\[2mm]}
    \onslide<+->{\D{F}{x\vphantom{y}}\D{x}{x} + \D{F}{y}\D{y}{x} + \D{F}{z}\D{z}{x} & = 0 \\[2mm]}
    \onslide<+->{\D{F}{x\vphantom{y}} + 0 + \D{F}{z}\D{z}{x}                        & = 0}
  \end{align*}
  \column{0.5\linewidth}
  \begin{align*}
    \onslide<+->{F_x + F_z \D{z}{x\vphantom{y}} & = 0                   \\[2mm]}
    \onslide<+->{F_z \D{z}{x\vphantom{y}}       & = - F_x               \\[2mm]}
    \onslide<+->{\D{z}{x\vphantom{y}}           & = - \frac{F_x}{F_z} \qquad \text{desde que } F_z \neq 0}
  \end{align*}
\end{columns}
\end{proofframe}

%------------------------------------------------------------------------------%
\begin{proofframe}{Derivação Implícita -- Justificativa}
  \begin{columns}
    \column{0.5\linewidth}
    \begin{align*}
      F\big(x, y, z(x, y)\big)                               & = 0 \\[2mm]
      \D{}{y}\left[F\big(x, y, z(x, y)\big)\right]           & = 0 \\[2mm]
      \D{F}{x}\D{x}{y} + \D{F}{y}\D{y}{y} + \D{F}{z}\D{z}{y} & = 0 \\[2mm]
      0 + \D{F}{y} + \D{F}{z}\D{z}{y}                        & = 0
    \end{align*}
    \column{0.5\linewidth}
    \begin{align*}
      F_y + F_z \D{z}{y} & = 0                   \\[2mm]
      F_z \D{z}{y}       & = - F_y               \\[2mm]
      \D{z}{y}           & = - \frac{F_y}{F_z} \qquad \text{desde que } F_z \neq 0
    \end{align*}
  \end{columns}
\end{proofframe}

%------------------------------------------------------------------------------%
\begin{frame}{Derivação Implícita -- Três Variáveis -- Fórmulas}
  Suponha que $F(x, y, z)$ seja diferenciável e que a equação
  \[
    F(x, y, z) = 0
  \]
  defina $z$ como função de $(x, y)$, então
  \[
    \D{z}{x} = -\frac{F_x}{F_z}
  \]
  \[
    \D{z}{y} = -\frac{F_y}{F_z}
  \]
  desde que \(F_z \neq 0\)
\end{frame}

%------------------------------------------------------------------------------%
\include{G-02-Derivacao_implicita-ex2}

%------------------------------------------------------------------------------%
\section{Exemplos}

\include{G-02-Derivacao_implicita-ex3}
\include{G-02-Derivacao_implicita-ex4}

%------------------------------------------------------------------------------%
\section{Lista Mínima}

%------------------------------------------------------------------------------%
\begin{frame}{Lista Mínima}
  Cálculo Vol. 2 do Thomas 12\textsuperscript{a} ed. -- Seção 14.4
  \begin{enumerate}
    \item Estudar o texto da seção
    \item Resolver os exercícios: 26, 28, 30, 32
  \end{enumerate}

  \vfill
  Atenção: A prova é baseada no livro, não nas apresentações
\end{frame}

%------------------------------------------------------------------------------%
\end{document}
%------------------------------------------------------------------------------%
